%%%%%%%%%%%%%%%%%%%%%%%%%%%%%%%%%%%%%%%%%%%%%%%%%%%%%%%%%%%%%%%%%%%%%
%% sec_theory_EM.tex
%%
%% Sections 2.2 -- 2.4 for paper2_main.tex
%% RMG Periodic RT-TDDFT -- Paper 2
%%
%% Covers:
%%   2.2  Interaction with electromagnetic radiation
%%   2.3  Gauge transformation: length to velocity gauge
%%   2.4  Vector potential and magnetic perturbations
%%
%% CONVENTIONS (all from notation.tex):
%%   Atomic units: hbar = m_e = e = 4*pi*eps_0 = 1
%%   Electron charge = -1  (e > 0 is elementary charge magnitude)
%%   Imaginary unit: \imath  (dotless i, consistent with paper 1)
%%   \myblue{} for all annotations and placeholders
%%   NEVER use \myred{}
%%
%% SOURCES USED:
%%   NEAT-notes-RMG_main.tex  lines 116-716  (primary)
%%   NEAT-project-notes_main.tex  lines 256-716  (supplementary)
%%   Discussion with Jacek, April 2026 (plane-wave/multipole framing)
%%%%%%%%%%%%%%%%%%%%%%%%%%%%%%%%%%%%%%%%%%%%%%%%%%%%%%%%%%%%%%%%%%%%%


%--------------------------------------------------------------------
\subsection{Interaction of Matter with Electromagnetic Radiation}
\label{sec:em-interaction}
%--------------------------------------------------------------------

The response of a periodic electronic system to an external
electromagnetic perturbation is the central quantity computed by
\texttt{RMG} in the time-dependent mode.
Before specializing to the particular perturbation scheme used in
this work, we briefly survey the general framework, which places
our implementation within a broader class of time-dependent
response calculations.

\subsubsection{The electromagnetic plane wave and minimal coupling}
\label{sec:plane-wave}

A monochromatic electromagnetic plane wave is described by the
vector potential
\begin{equation}
  \Avec(\mathbf{r},t)
  = A_0\,\hat{\varepsilon}\,
    e^{\imath(\mathbf{q}\cdot\mathbf{r}-\omega t)}
    + \text{c.c.},
  \label{eq:plane-wave-A}
\end{equation}
where $\hat{\varepsilon}$ is the polarization unit vector,
$\mathbf{q}$ is the photon wavevector with $|\mathbf{q}|=\omega/c$,
and $\omega$ is the angular frequency.
In the Coulomb gauge ($\nabla\cdot\Avec=0$, $\phi=0$) the transversality
condition requires $\hat{\varepsilon}\cdot\mathbf{q}=0$: the field
oscillates perpendicular to the direction of propagation.
The physical electric and magnetic fields follow from the potentials:
\begin{align}
  \Evec(\mathbf{r},t)
  &= -\frac{\partial\Avec}{\partial t}
   = \imath\omega A_0\,\hat{\varepsilon}\,
     e^{\imath(\mathbf{q}\cdot\mathbf{r}-\omega t)}
     + \text{c.c.},
  \label{eq:plane-wave-E}
  \\
  \Bvec(\mathbf{r},t)
  &= \nabla\times\Avec
   = \imath\mathbf{q}\times
     \!\left(A_0\hat{\varepsilon}\right)
     e^{\imath(\mathbf{q}\cdot\mathbf{r}-\omega t)}
     + \text{c.c.},
  \label{eq:plane-wave-B}
\end{align}
where in eq.~(\ref{eq:plane-wave-B}) we used
$\nabla e^{\imath\mathbf{q}\cdot\mathbf{r}}
= \imath\mathbf{q}\,e^{\imath\mathbf{q}\cdot\mathbf{r}}$.
Since $|\mathbf{q}|=\omega/c$, the magnitudes satisfy
$|\Bvec|/|\Evec| = c|\mathbf{q}|/\omega = 1$ in atomic units with
$c=1/\alpha$, confirming the equal-magnitude relation for a free
plane wave.

The coupling of the electromagnetic field to a Kohn-Sham electron
is introduced through minimal substitution
$\hat{\mathbf{p}} \to \hat{\mathbf{p}} + \Avec(\mathbf{r},t)$
(atomic units, electron charge $-1$),
giving the time-dependent Kohn-Sham Hamiltonian
\begin{equation}
  \hat{H}(t)
  = \frac{1}{2}
    \!\left(\hat{\mathbf{p}} + \Avec(\mathbf{r},t)\right)^{\!2}
    + V_\mathrm{KS}(\mathbf{r},t).
  \label{eq:H-minimal}
\end{equation}
Expanding the kinetic term and retaining only the linear (first-order)
coupling in $\Avec$ (valid in the linear response regime where
$|A_0|\ll 1$ in atomic units):
\begin{equation}
  \hat{H}(t)
  = \hat{H}_0
  + \underbrace{\Avec(\mathbf{r},t)\cdot\hat{\mathbf{p}}}_{\hat{H}_\mathrm{int}(t)}
  + \mathcal{O}(A_0^2),
  \label{eq:H-linear}
\end{equation}
where $\hat{H}_0 = -\nabla^2/2 + V_\mathrm{KS}$ is the unperturbed
Kohn-Sham Hamiltonian and the diamagnetic $\Avec^2$ term is dropped
as it contributes at second order in the applied
field.\cite{yabana-bertsch-1996,Mattiat-Luber-2022}
The interaction Hamiltonian is therefore
\begin{equation}
  \hat{H}_\mathrm{int}(\mathbf{r},t)
  = A_0\,\hat{\varepsilon}\cdot\hat{\mathbf{p}}\;
    e^{\imath(\mathbf{q}\cdot\mathbf{r}-\omega t)}
    + \text{c.c.}
  \label{eq:Hint-full}
\end{equation}

\subsubsection{The multipole expansion and the long-wavelength approximation}
\label{sec:multipole}

The spatial dependence of the interaction~(\ref{eq:Hint-full}) enters
entirely through the phase factor $e^{\imath\mathbf{q}\cdot\mathbf{r}}$.
Expanding in a Taylor series about $\mathbf{q}\cdot\mathbf{r}=0$:
\begin{equation}
  e^{\imath\mathbf{q}\cdot\mathbf{r}}
  = 1
  + \imath(\mathbf{q}\cdot\hat{\mathbf{r}})
  + \frac{1}{2}\!\left[\imath(\mathbf{q}\cdot\hat{\mathbf{r}})\right]^2
  + \cdots
  \label{eq:multipole-expansion}
\end{equation}
Inserting eq.~(\ref{eq:multipole-expansion}) into the matrix element
of $\hat{H}_\mathrm{int}$ between Kohn-Sham orbitals $|\psi_a\rangle$
and $|\psi_b\rangle$:
\begin{align}
  \langle\psi_b|\hat{H}_\mathrm{int}|\psi_a\rangle
  &\propto
  \underbrace{
    \langle\psi_b|\hat{\varepsilon}\cdot\hat{\mathbf{p}}|\psi_a\rangle
  }_{\text{E1: electric dipole (velocity form)}}
  +\;\imath\!
  \underbrace{
    \langle\psi_b|(\mathbf{q}\cdot\hat{\mathbf{r}})
    (\hat{\varepsilon}\cdot\hat{\mathbf{p}})|\psi_a\rangle
  }_{\text{M1 + E2: magnetic dipole + electric quadrupole}}
  + \cdots
  \label{eq:multipole-me}
\end{align}
The successive terms are the electric dipole (E1), magnetic dipole
plus electric quadrupole (M1+E2), and higher multipole contributions.
Each term carries an additional factor of $|\mathbf{q}|\,a_0$ relative
to the preceding one.
For optical and ultraviolet radiation, $|\mathbf{q}| = \omega/c$;
at $\lambda = 400$~nm (near UV), $|\mathbf{q}| \approx 8\times10^{-4}$~a$_0^{-1}$,
while a unit cell of linear dimension $a \sim 10$~\AA\ gives
$|\mathbf{q}|\,a \approx 1.5\times10^{-2}$.
The multipole series therefore converges rapidly and the
\emph{long-wavelength} (electric dipole) \emph{approximation}
\begin{equation}
  e^{\imath\mathbf{q}\cdot\mathbf{r}} \approx 1,
  \qquad |\mathbf{q}|\,a \ll 1,
  \label{eq:LWA}
\end{equation}
is quantitatively excellent for all optical and UV spectroscopy.
Under this approximation the vector potential becomes spatially
uniform: $\Avec(\mathbf{r},t)\to\Avec(t)$, and the interaction
Hamiltonian reduces to
\begin{equation}
  \hat{H}_\mathrm{int}(t)
  = \Avec(t)\cdot\hat{\mathbf{p}}.
  \label{eq:Hint-LWA}
\end{equation}

\subsubsection{Crystal momentum conservation: vertical transitions}
\label{sec:vertical}

For a periodic crystal with Bloch states
$|\psi_{n\mathbf{k}}\rangle$, the matrix element of the
long-wavelength interaction~(\ref{eq:Hint-LWA}) between states at
crystal momenta $\mathbf{k}$ and $\mathbf{k}'$ involves the integral
of $e^{\imath(\mathbf{k}-\mathbf{k}')\cdot\mathbf{r}}$ over the unit
cell, which vanishes unless $\mathbf{k}=\mathbf{k}'$:
\begin{equation}
  \langle\psi_{m\mathbf{k}'}|
  \Avec(t)\cdot\hat{\mathbf{p}}
  |\psi_{n\mathbf{k}}\rangle
  \propto \delta_{\mathbf{k},\mathbf{k}'}.
  \label{eq:vertical}
\end{equation}
The uniform field $\Avec(t)$ therefore connects only states at the
\emph{same} crystal momentum: all induced transitions are
\emph{vertical} ($\Delta\mathbf{k}=0$) in the Brillouin zone.

More generally, the full interaction~(\ref{eq:Hint-full}) with the
phase factor $e^{\imath\mathbf{q}\cdot\mathbf{r}}$ retained would
give $\delta_{\mathbf{k}',\mathbf{k}+\mathbf{q}}$: the photon
deposits its wavevector $\mathbf{q}$ into the electronic system,
shifting the crystal momentum by $\Delta\mathbf{k}=\mathbf{q}$.
For optical photons this shift is negligible ($|\mathbf{q}|\ll\pi/a$),
but the same formal machinery, applied with $\mathbf{q}$ as a free
parameter decoupled from $\omega$, forms the theoretical basis for
computing finite-wavevector response functions
$\chi(\mathbf{q},\omega)$ directly from RT-TDDFT.\cite{yabana-bertsch-1996}
In that framework the electronic system is perturbed by a
density-wave kick of the form
$\delta V(\mathbf{r},t)=V_0\,e^{\imath\mathbf{q}\cdot\mathbf{r}}\delta(t)$,
and the Fourier component of the density response $\delta\rho(\mathbf{q},t)$
yields $\chi(\mathbf{q},\omega)$ and thence the dynamic structure
factor $S(\mathbf{q},\omega) \propto \mathrm{Im}\,\chi(\mathbf{q},\omega)$
after Fourier transformation in time.
Repeating the calculation for a sequence of commensurate wavevectors
$\mathbf{q}$ gives access to plasmon dispersion
curves,\cite{yabana-bertsch-1996} interband transitions at finite
momentum transfer, and quantities measurable by inelastic X-ray
scattering or electron energy loss
spectroscopy.\myblue{[cite: EELS RT-TDDFT ref; IXS ref]}
This extension of the present implementation is deferred to future
work; in this paper we restrict to the optical ($\mathbf{q}\to 0$)
regime throughout.

\myblue{[Note: An alternative perturbation that gives spatially
resolved response -- analogous to a focused electron beam -- is the
Coulomb potential of a point charge placed at position $\mathbf{R}$:
$\delta V(\mathbf{r}) = V_0/|\mathbf{R}-\mathbf{r}|$. This approach
was implemented and validated in ref.~\cite{jakowski-tddft-coulomb-2019}
\myblue{[check cite key]}. The Coulomb perturbation is the real-space
counterpart of the plane-wave perturbation: by the Fourier expansion
$1/|\mathbf{R}-\mathbf{r}| = (4\pi/\Omega)\sum_\mathbf{q}
e^{\imath\mathbf{q}\cdot(\mathbf{r}-\mathbf{R})}/|\mathbf{q}|^2$,
it is a superposition of all $e^{\imath\mathbf{q}\cdot\mathbf{r}}$
modes weighted by $1/|\mathbf{q}|^2$, and thus samples the full
$\chi(\mathbf{q},\omega)$ at all wavevectors simultaneously but
with a $1/q^2$ weight that emphasizes long-wavelength
components.]}


%--------------------------------------------------------------------
\subsection{Gauge Transformation: Length to Velocity Gauge}
\label{sec:gauge}
%--------------------------------------------------------------------

\subsubsection{Why the length gauge fails for periodic systems}
\label{sec:LG-failure}

In the long-wavelength limit the two standard gauge choices are:
\begin{alignat}{3}
  &\text{Length gauge (LG):}\quad
  &&\Avec^\mathrm{LG} = \mathbf{0},
  \quad
  &&\phi^\mathrm{LG}(\mathbf{r},t)
    = -\hat{\mathbf{r}}\cdot\mathbf{E}(t),
  \label{eq:LG-potentials}
  \\
  &\text{Velocity gauge (VG):}\quad
  &&\Avec^\mathrm{VG}(t) = -\!\int_{-\infty}^t\!\!\mathbf{E}(t')\,dt',
  \quad
  &&\phi^\mathrm{VG} = 0.
  \label{eq:VG-potentials}
\end{alignat}
In the length gauge the interaction Hamiltonian takes the familiar
electric dipole form $\hat{H}_\mathrm{int}^\mathrm{LG}
= \hat{\mathbf{r}}\cdot\mathbf{E}(t)$,
which requires evaluating matrix elements of the position operator
$\langle\psi_a|\hat{r}_\alpha|\psi_b\rangle$ between the propagated
orbitals.

For a finite molecular system, as in paper~1,\cite{jakowski-rmg-tddft-2025}
these matrix elements are well-defined and
the length gauge is the natural choice.
For a periodic system with Born-von Kármán boundary conditions, however,
the position operator $\hat{\mathbf{r}}$ is incompatible with the
periodicity of the Bloch states: translating by a lattice vector
$\mathbf{R}$ maps $\hat{r}_\alpha$ to $\hat{r}_\alpha + R_\alpha$,
changing the operator and making its expectation values
ill-defined.\cite{resta-1998,king-smith-vanderbilt-1993}
Equivalently, the matrix element
$\langle\psi_{n\mathbf{k}}|\hat{r}_\alpha|\psi_{m\mathbf{k}}\rangle$
diverges for extended Bloch states.
The velocity gauge resolves this by removing $\hat{\mathbf{r}}$ from
the interaction entirely and replacing it with the momentum operator
$\hat{\mathbf{p}}$, whose matrix elements between Bloch states are
well-defined (Section~\ref{sec:current}).

\subsubsection{The Göppert-Mayer gauge transformation}
\label{sec:GM-transform}

The length and velocity gauges are related by the gauge function
\begin{equation}
  \chi(\mathbf{r},t) = -\hat{\mathbf{r}}\cdot\Avec^\mathrm{VG}(t),
  \label{eq:gauge-function}
\end{equation}
which generates the Göppert-Mayer transformation.\cite{goeppert-mayer-1931}
Under a gauge transformation with function $\chi$, the potentials,
wavefunctions, and operators transform as
\begin{align}
  \Avec &\to \Avec + \nabla\chi,
  \label{eq:gauge-A}
  \\
  \phi &\to \phi - \frac{\partial\chi}{\partial t},
  \label{eq:gauge-phi}
  \\
  \Psi(\mathbf{r},t) &\to \Psi'(\mathbf{r},t)
  = e^{-\imath\chi(\mathbf{r},t)}\Psi(\mathbf{r},t)
  = \hat{U}(\mathbf{r},t)\Psi(\mathbf{r},t),
  \label{eq:gauge-psi}
  \\
  \hat{O} &\to \hat{O}' = \hat{U}\hat{O}\hat{U}^\dagger,
  \label{eq:gauge-O}
\end{align}
where $\hat{U}(\mathbf{r},t) = e^{-\imath\chi(\mathbf{r},t)}$
is unitary (so charge density
$\rho'=|\Psi'|^2=|\Psi|^2=\rho$ is unchanged)
and in atomic units with electron charge $-1$ the phase is
$\chi = -\hat{\mathbf{r}}\cdot\Avec(t)$ giving
$\hat{U} = e^{\imath\hat{\mathbf{r}}\cdot\Avec(t)}$.
Applying eqs.~(\ref{eq:gauge-A})--(\ref{eq:gauge-phi}) to the
LG potentials with $\chi = -\hat{\mathbf{r}}\cdot\Avec^\mathrm{VG}$
one verifies that $\Avec^\mathrm{LG}+\nabla\chi = \Avec^\mathrm{VG}$
and $\phi^\mathrm{LG} - \partial_t\chi = \phi^\mathrm{VG} = 0$,
confirming that eq.~(\ref{eq:gauge-function}) is indeed the
transformation connecting the two gauges.

\subsubsection{The velocity-gauge Kohn-Sham Hamiltonian}
\label{sec:H-VG}

The gauge-transformed Hamiltonian is obtained from
$\hat{H}' = \hat{U}\hat{H}\hat{U}^\dagger
- \imath\hat{U}\partial_t\hat{U}^\dagger$.
We evaluate the two contributions in turn.

\paragraph{Kinetic energy.}
The transformation of the canonical momentum operator under $\hat{U}$
follows from evaluating $\hat{U}\hat{p}_\alpha\hat{U}^\dagger$
on an arbitrary function $g(\mathbf{r})$:
\begin{equation}
  \hat{U}\hat{p}_\alpha\hat{U}^\dagger g
  = e^{\imath\mathbf{r}\cdot\Avec}\,(-\imath\nabla_\alpha)
    \!\left(e^{-\imath\mathbf{r}\cdot\Avec}g\right)
  = \hat{p}_\alpha g - A_\alpha g,
  \label{eq:U-p-Udag}
\end{equation}
where we used $\nabla_\alpha e^{-\imath\mathbf{r}\cdot\Avec}
= -\imath A_\alpha e^{-\imath\mathbf{r}\cdot\Avec}$.
Therefore $\hat{U}\hat{p}_\alpha\hat{U}^\dagger = \hat{p}_\alpha - A_\alpha$,
and the kinetic term transforms as
\begin{equation}
  \hat{U}\frac{\hat{p}^2}{2}\hat{U}^\dagger
  = \frac{1}{2}\!\left(\hat{\mathbf{p}} - \Avec\right)^{\!2}.
  \label{eq:U-T-Udag}
\end{equation}
Since the LG Hamiltonian for an electron (charge $-1$) has kinetic
term $\frac{1}{2}\hat{p}^2$ and the gauge-transformed kinetic
momentum is $\hat{\mathbf{p}} - \Avec$, the electron kinetic energy
in the velocity gauge becomes $\frac{1}{2}(\hat{\mathbf{p}} + \Avec)^2$
after accounting for the electron charge sign
(the sign reversal $-\Avec\to+\Avec$ comes from the factor of $-1$
in the electron charge).\footnote{%
In Gaussian units with electron charge $q=-e$, the minimal substitution
gives $\hat{\mathbf{p}}\to\hat{\mathbf{p}}-q\Avec/c
=\hat{\mathbf{p}}+e\Avec/c$; in atomic units $e=1$, $c=1/\alpha$ is
absorbed into $\Avec$, giving $\hat{\mathbf{p}}+\Avec$ as written.}

\paragraph{Time-derivative term.}
The second contribution is
\begin{equation}
  -\imath\hat{U}\frac{\partial\hat{U}^\dagger}{\partial t}
  = -\imath\,e^{\imath\mathbf{r}\cdot\Avec}
    \frac{\partial}{\partial t}\!
    \left(e^{-\imath\mathbf{r}\cdot\Avec}\right)
  = -\mathbf{r}\cdot\dot{\Avec}(t)
  = \mathbf{r}\cdot\mathbf{E}(t),
  \label{eq:U-dt-Udag}
\end{equation}
where we used $\mathbf{E}(t) = -\partial\Avec/\partial t$
(Coulomb gauge, atomic units).
This term exactly cancels the length-gauge coupling
$\mathbf{r}\cdot\mathbf{E}(t)$, as required.

\paragraph{Local potential.}
Since $V_\mathrm{loc}(\mathbf{r})$ is diagonal in position space,
$\hat{U}V_\mathrm{loc}\hat{U}^\dagger = V_\mathrm{loc}$:
the local potential is gauge invariant.

\paragraph{Nonlocal pseudopotential.}
The nonlocal pseudopotential $\Vnl$ is not diagonal in position
space; its gauge transformation is
\begin{equation}
  \hat{U}\Vnl\hat{U}^\dagger
  = e^{\imath\hat{\mathbf{r}}\cdot\Avec}\Vnl
    e^{-\imath\hat{\mathbf{r}}\cdot\Avec}
  = \Vnl + \imath\Avec(t)\cdot[\hat{\mathbf{r}},\Vnl]
    + \mathcal{O}(A_0^2).
  \label{eq:U-Vnl-Udag}
\end{equation}
The first-order correction $\imath\Avec\cdot[\hat{\mathbf{r}},\Vnl]$
is the gauge-transformation contribution to the nonlocal
pseudopotential in the velocity gauge.
As shown in Appendix~\ref{app:gauge_nlpp}, this correction is
equivalent to the Starace commutator term that appears in the
velocity-gauge current operator (Section~\ref{sec:current}):
the two derivations --- gauge transformation of $\Vnl$ and
Heisenberg equation of motion for $\hat{v}_\alpha$ --- yield the
same correction, confirming internal
consistency.\cite{starace-1971,Mattiat-Luber-2022}

\paragraph{Full velocity-gauge TDKS Hamiltonian.}
Collecting all contributions, the length-gauge TDKS Hamiltonian
transforms to the velocity-gauge form
\begin{equation}
  \boxed{
  \hat{H}^\mathrm{VG}(t)
  = \frac{1}{2}\!\left(\hat{\mathbf{p}} + \Avec(t)\right)^{\!2}
  + V_\mathrm{KS}(\mathbf{r},t)
  + \Vnl
  + \imath\Avec(t)\cdot\left[\hat{\mathbf{r}},\Vnl\right],
  }
  \label{eq:H-VG}
\end{equation}
where $V_\mathrm{KS}$ includes the local external, Hartree, and
exchange-correlation potentials, and the last term is the
first-order nonlocal pseudopotential correction from the gauge
transformation.
Equation~(\ref{eq:H-VG}) is the velocity-gauge TDKS Hamiltonian
used in \texttt{RMG} for all periodic RT-TDDFT calculations in this
work.\cite{Mattiat-Luber-2022,pickard-mauri-2001}

For the periodic case, the Bloch ansatz
$\psi_{n\mathbf{k}}(\mathbf{r},t)
= e^{\imath\mathbf{k}\cdot\mathbf{r}}u_{n\mathbf{k}}(\mathbf{r},t)$
is substituted into eq.~(\ref{eq:H-VG}); the Bloch factor
contributes an additional $+\mathbf{k}$ shift to the canonical
momentum (see Section~\ref{sec:current-bloch}), giving the
$\mathbf{k}$-dependent Hamiltonian for the cell-periodic part
$u_{n\mathbf{k}}$:
\begin{equation}
  \hat{H}^\mathrm{VG}_\mathbf{k}(t)
  = \frac{1}{2}\!\left(-\imath\nabla + \mathbf{k} + \Avec(t)\right)^{\!2}
  + V_\mathrm{KS}(\mathbf{r},t)
  + \Vnl
  + \imath\Avec(t)\cdot\left[\hat{\mathbf{r}},\Vnl\right].
  \label{eq:H-VG-k}
\end{equation}

\subsubsection{The perturbative momentum kick}
\label{sec:delta-kick}

In the linear response regime the vector potential is applied as a
single instantaneous impulse at $t=0$:
\begin{equation}
  \Avec(t) = A_0\,\hat{\varepsilon}\,\delta(t),
  \label{eq:delta-kick}
\end{equation}
which corresponds to a constant electric field pulse of infinitesimal
duration and amplitude $E_0 = A_0$ in the frequency
domain.\cite{yabana-bertsch-1996}
After the kick ($t>0$), $\Avec(t)=0$ and the Hamiltonian~(\ref{eq:H-VG})
reduces to $\hat{H}_0 + V_\mathrm{KS}[\rho(t)]$, with the time
evolution driven entirely by the self-consistent response of the
density.
The delta kick imparts a phase $e^{\imath A_0\hat{\varepsilon}\cdot\hat{\mathbf{p}}}$
to each orbital at $t=0^+$, which in the small-amplitude limit
$|A_0|\ll 1$ reduces to a uniform momentum shift:
$\psi_{n\mathbf{k}}(\mathbf{r},0^+)
\approx (1 + \imath A_0\hat{\varepsilon}\cdot\hat{\mathbf{p}})
\psi_{n\mathbf{k}}(\mathbf{r},0)$.
This perturbation is identical for all $\mathbf{k}$-points and all
bands, excites all frequencies simultaneously, and the full optical
response tensor is recovered from three independent simulations with
kick directions $\hat{\varepsilon}\in\{\hat{x},\hat{y},\hat{z}\}$.


%--------------------------------------------------------------------
\subsection{Vector Potential and Magnetic Perturbations}
\label{sec:vector-potential}
%--------------------------------------------------------------------

The velocity-gauge Hamiltonian~(\ref{eq:H-VG}) provides a unified
framework for both electric and magnetic perturbations, distinguished
solely by the spatial structure of the vector potential $\Avec(\mathbf{r},t)$.

\subsubsection{Electric perturbation: spatially uniform A(t)}
\label{sec:electric-pert}

In the long-wavelength approximation (Section~\ref{sec:multipole}),
$\Avec(\mathbf{r},t)\to\Avec(t)$ is spatially uniform.
The magnetic field associated with a spatially uniform vector potential
vanishes identically:
\begin{equation}
  \Bvec(t) = \nabla\times\Avec(t) = \mathbf{0},
  \label{eq:B-uniform-zero}
\end{equation}
so the perturbation is purely electric.
This is the regime of the present work: the delta-kick
implementation (Section~\ref{sec:delta-kick}) applies a uniform
$\Avec(t)$ and measures the resulting current response, from which
the optical conductivity is extracted (Section~\ref{sec:postproc}).

\subsubsection{Magnetic perturbation: spatially varying A(r)}
\label{sec:magnetic-pert}

A static or time-dependent magnetic field $\Bvec$ corresponds to a
spatially varying vector potential.
In the symmetric gauge centered at origin $\mathbf{R}_G$:
\begin{equation}
  \Avec(\mathbf{r})
  = \tfrac{1}{2}\Bvec\times(\mathbf{r}-\mathbf{R}_G),
  \label{eq:symmetric-gauge}
\end{equation}
from which $\nabla\times\Avec = \Bvec$ follows directly.
For a time-dependent $\Bvec(t)$, the orbital magnetic coupling
$\Avec(\mathbf{r},t)\cdot\hat{\mathbf{p}}$ in eq.~(\ref{eq:H-VG})
is no longer spatially uniform and drives transitions with
$\Delta\mathbf{k}\neq 0$ within the unit cell (for $|\mathbf{q}|
= |\Bvec|/|\Avec| \cdot a \ll 1$, still small but in principle
non-zero).
This framework is the foundation for computing magnetically induced
circular dichroism (ECD/MCD) spectra from first
principles.\cite{Mattiat-Luber-2022}

A subtlety of eq.~(\ref{eq:symmetric-gauge}) is the dependence on the
gauge origin $\mathbf{R}_G$.
For a complete basis set, physical observables are independent of
this choice; for any finite basis, including the real-space grid used
in \texttt{RMG}, residual gauge-origin dependence appears and vanishes
only in the limit of fine grid spacing.
On a sufficiently dense grid --- which is already required for
accurate DFT total energies --- this dependence is negligible in
practice.\myblue{[cite Mattiat for similar discussion in Gaussian
basis context]}

\subsubsection{Scope of the present implementation}
\label{sec:scope}

The calculations presented in this paper employ the electric
perturbation scheme of Section~\ref{sec:electric-pert} exclusively.
The spatially uniform delta kick~(\ref{eq:delta-kick}) gives
$\Bvec=0$ for all $t>0$, and the diamagnetic term $\Avec^2$ likewise
vanishes (Section~\ref{sec:para-dia}).
The extension to magnetic perturbations~(\ref{eq:symmetric-gauge}),
which would enable ECD and MCD spectroscopy, is a natural next step
within the same velocity-gauge framework and is left for future
work.
